\chapter{Introduction}

\section{Background}

Large language model (LLM) agents are increasingly designed as systems that combine reasoning, action, memory, and external capability use rather than as standalone text generators. Early work on retrieval-augmented generation showed that language models can improve performance by accessing external information at inference time rather than relying solely on knowledge stored in model parameters \cite{lewis2021rag, izacard2022atlas}. Later work on tool use extended this idea from document retrieval to action execution by enabling models to call external functions and APIs \cite{schick2023toolformer, patil2023gorilla, qin2023toolllm}.

More recently, agent systems have begun to treat reusable capabilities as \emph{skills}. A skill is not simply an atomic tool call. It may include instructions, workflow steps, scripts, examples, dependencies, input requirements, output expectations, and constraints. In practical systems such as filesystem-based agent skill libraries, the full skill artifact may be stored as a directory containing a \texttt{SKILL.md} file plus optional resources. The agent may initially see only a compressed representation, such as a name and description, before deciding whether to load the full artifact.

This shift changes the retrieval problem. If the object being selected is a high-level procedure rather than an atomic function, retrieval quality depends not only on topical relevance, but also on procedural suitability. Two skills may look semantically similar because they refer to the same domain or task family, while still differing in preconditions, workflow steps, required resources, output format, or side effects.

\section{Problem Statement}

The problem addressed in this thesis is semantic confusion in large agent skill libraries. As skill libraries grow, a user request may match many plausible skills at the level of names and descriptions. The correct skill may exist, but the retrieval system may surface a nearby skill that sounds relevant while being procedurally wrong.

This failure mode is different from ordinary retrieval failure. The wrong result is often not unrelated. It is close enough to be tempting. For example, a request to ``extract information from this PDF'' may plausibly match document review, document summarisation, citation grounding, field extraction, or layout-preserving conversion. These skills may all mention documents, extraction, or PDFs, but they require different procedures and produce different outputs. A flat description may not preserve the procedural distinctions required for correct selection.

The scalability issue makes this worse. In a small library, the main agent may be able to inspect all skill names and short descriptions in context. In a larger library, even this metadata can become too large or noisy. The system then needs a candidate-subsetting stage before the main agent loads any full skill document. This makes the skill representation critical: the correct skill must be surfaced before the full artifact can be inspected.

\section{Research Questions}

This thesis uses the following working research questions:

\textbf{RQ1:} Which operational information types help distinguish semantically similar but procedurally distinct agent skills?

\textbf{RQ2:} How do skill representation and retrieval-pipeline choices affect the preservation and use of the routing-relevant operational information identified in RQ1 under semantic confusability, and what accuracy, candidate-recall, and retrieval-cost trade-offs result?

These questions deliberately narrow the thesis from a broad comparison of ``flat'' and ``structured'' skill representation and retrieval to a diagnostic study of routing-relevant information. RQ1 asks which kinds of operational information actually make near-neighbour skills distinguishable. RQ2 separates a matched-content mechanism study from a full-library validation stage. RQ2a fixes the query, three sibling candidates, reviewed operational propositions, and selector while varying how the same facts are organised, flattened, expressed as prose, diluted, or explicitly scored field by field. RQ2b then evaluates four source-grounded representations---I1 native name--description metadata, I2 complete skill bodies, I3-flat extracted evidence, and field-labelled I3C---over a 2,433-skill library and 381 strict-gold prompts. It separates three first-stage retrievers from two fixed-Top-20 rerankers, so candidate inclusion, final ranking, and measured selector cost remain distinguishable.

This framing also limits the novelty claim. The thesis does not propose a new universal schema for agent skills. Recent structured skill representations such as Scheduling-Structural-Logical (SSL) already decompose skill artifacts into invocation-level signals, execution structure, and action/resource evidence. The contribution here is diagnostic and empirical: given semantically similar candidate skills, it tests which operational fields make the correct skill distinguishable and which skill retrieval methods can use those fields reliably.

\section{Scope}

The thesis focuses on retrieval and candidate selection before full skill execution. It does not attempt to build a complete general-purpose agent operating system. The main agent, the retriever, and the reranker are treated as separable components. The retriever may be a lexical method, an embedding model, a deterministic schema scorer, a neural reranker, or a hybrid pipeline. What matters is whether the retrieval stage can produce a candidate set that contains and ranks the procedurally correct skill.

The evaluation focuses on skill libraries where some skills are intentionally confusable. This is important because broad relevance retrieval is too easy for the research question. A benchmark that only asks whether a document task retrieves a document skill would not test the central failure mode. The difficult case is selecting between skills that share the same broad domain but differ in operational requirements.

\section{Contributions}

This thesis makes five contributions:

\begin{enumerate}
  \item A field-isolated benchmark design for skill retrieval under semantic confusability and procedural distinctness, with controlled near-neighbour clusters that vary one operational signal at a time.
  \item Empirical evidence about which operational information types act as routing signals, distinguishing strong positive routing signals, conditional or tie-breaker signals, and information that mainly supports execution after selection.
  \item A matched-content RQ2a comparison that separates operational-fact content, explicit field organisation, prose form, information dilution, and selector behaviour under lexical, single-vector semantic, cross-encoder, and semantic field-aware selection.
  \item A source-grounded RQ2b full-library validation over 2,433 skills and 381 strict-gold prompts, comparing I1, I2, I3-flat, and I3C under lexical, Qwen, and SkillRouter first-stage retrieval plus two fixed-Top-20 rerankers; it reports an accuracy--compression trade-off rather than a universal I3C gain.
  \item A failure-mode and efficiency analysis of when routing-relevant information is absent, diluted, implicit, misleading, or present but not exploited, while treating tree/graph architectures and downstream execution as future work.
\end{enumerate}

The contribution is therefore positioned as a controlled analysis of discriminative skill information, not as the first proposal of structured skill representation.

\section{Thesis Structure}

Chapter 2 reviews the literature on retrieval-augmented language models, tool use, agent skills, skill-library scaling, and structure-aware retrieval. Chapter 3 defines the vocabulary and layer distinctions used throughout the thesis. Chapter 4 describes the benchmark design. Chapter 5 explains the methodology and evaluation metrics. Chapter 6 reports the results. Chapter 7 discusses implications and failure modes. Chapter 8 concludes the thesis and outlines future work.
